\documentclass[11pt, a4paper]{article}

\usepackage[T1]{fontenc}
\usepackage[utf8]{inputenc}
\usepackage{tgpagella}
\usepackage[spanish, es-noshorthands]{babel}
\usepackage{amsmath, amssymb}
\usepackage[most]{tcolorbox}
\usepackage{listings}
\usepackage[margin=2.5cm]{geometry}
\usepackage{fancyhdr}

\pagestyle{fancy}
\fancyhf{}
\setlength{\headheight}{16pt}
\lhead{\textbf{UCB Sede Tarija} | Ing. Mecatrónica}
\rhead{IMT-221 | Trabajo Asincrónico S06-2}
\renewcommand{\headrulewidth}{0.4pt}

\definecolor{ucbblue}{RGB}{0, 51, 102}

\lstset{
    language=Python,
    backgroundcolor=\color{gray!10},
    basicstyle=\ttfamily\small,
    keywordstyle=\color{blue},
    commentstyle=\color{green!50!black}\itshape,
    stringstyle=\color{red},
    frame=single,
    breaklines=true,
    numbers=left,
    numberstyle=\tiny\color{gray}
}

\begin{document}

\section*{Consolidación Asincrónica: Dependencia de Parámetros}

\subsection*{Parte A y C: Ejercicios Numéricos}
\textbf{A.} Asume $I_{CQ}=1.2\,\text{mA}$, $\beta=120$, $T \approx 300\,\text{K}$. Calcula los valores de $g_m, r_e, y r_\pi$. \\
\textbf{C.} Para la etapa de emisor común con $R_C=2.7\,\text{k}\Omega$ y $R_L=10\,\text{k}\Omega$, calcula la ganancia de voltaje $A_v$. (Asume emisor en tierra AC y $r_o \to \infty$).

\subsection*{Parte B y D: Evaluación Conceptual}
Responde en tu cuaderno:
\begin{enumerate}
    \item ¿Por qué $g_m$ depende del punto de operación $Q$?
    \item Al transformar un circuito a su equivalente AC, ¿por qué la fuente $V_{CC}$ se convierte en tierra pero el resistor $R_C$ no desaparece?
    \item Dibuja de memoria el modelo Híbrido-$\pi$ (incluyendo nodos, resistencias y la dirección de la fuente de corriente).
    \item ¿Por qué la fórmula $-g_m R_C$ puede fallar al predecir la salida si la señal de entrada es demasiado grande?
\end{enumerate}

\subsection*{Tabla de Parámetros en Python (Opcional)}
Analiza este código. ¿Cuál es la tendencia de los parámetros cuando $I_C$ aumenta?

\begin{lstlisting}
import numpy as np

VT = 25.85e-3      # V at approximately 300 K
beta = 100.0

IC = np.array([0.5e-3, 1e-3, 2e-3, 4e-3])

gm = IC / VT
re = 1.0 / gm
rpi = beta / gm

for ic, gm_i, re_i, rpi_i in zip(IC, gm, re, rpi):
    print(
        f"IC={ic*1e3:4.1f} mA | "
        f"gm={gm_i*1e3:6.2f} mS | "
        f"re={re_i:7.2f} ohm | "
        f"rpi={rpi_i/1e3:6.2f} kohm"
    )
\end{lstlisting}

\end{document}
